% GENERATED FILE. Edit canonical lessons, then run npm run generate:books.
% canonical-source-hash: fnv1a64:9801a0d6d2283b9c
% canonical-lessons: GU-C210-erport, GU-C210-pasport, GU-C210-parvano, GU-C210-ladvu, GU-C210-uncakvu

\chapter{Papers and baggage}
\label{ch:gu-papers-and-baggage}
\hlchaptermodality{210}

\begin{chapteropening}
\textbf{By the end of this chapter:} \emph{I can say what we showed and carried at the airport.}

Complete the last lesson of chapter 210: I can say what we showed and carried at the airport.
\end{chapteropening}

\section[\texorpdfstring{\gu{એરપોર્ટ} (erporṭ)}{એરપોર્ટ (erporṭ)}]{\gu{એરપોર્ટ} (erporṭ) — airport}

\label{lesson:GU-C210-erport}



\begin{quote}
Before the new one: say the Gujarati for to get stuck, then the Gujarati for to slide.
\end{quote}

\subsection*{You'll want to know: \gu{એરપોર્ટ}}
\textbf{\gu{એરપોર્ટ}} — \emph{erporṭ} — \textquotedblleft{}airport\textquotedblright{}.

\begin{quote}
\textbf{\gu{અમે} \gu{બપોરે} \gu{એરપોર્ટ} \gu{ગયા}.} — \emph{ame bapore erporṭ gayā} — \textquotedblleft{}We went to the airport at midday.\textquotedblright{}
\end{quote}

\begin{cousinweb}
\textbf{\gu{એરપોર્ટ}} (\emph{erporṭ}) is a loan from English \emph{airport}.
\end{cousinweb}

\begin{grammarlens}[title={What you've built}]
One more word for everyday life at A2.
\end{grammarlens}

\subsection*{Your turn}
\begin{itemize}
\raggedright
  \item \emph{Say it:} \emph{erporṭ}
  \item \emph{Say it:} \emph{erporṭ}, once more
  \item \emph{Say it:} \emph{ame bapore erporṭ gayā} — We went to the airport at midday.
  \item \emph{Recall:} say the Gujarati for to get stuck, then the Gujarati for to slide, then say \emph{erporṭ} again
\end{itemize}

\begin{tcolorbox}[breakable,colback=teal!4,colframe=teal!35!black,title={Before you move on}]
What does \gu{એરપોર્ટ} mean? (\textquotedblleft{}Airport\textquotedblright{}.) Say it once more.
\end{tcolorbox}
\section[\texorpdfstring{\gu{પાસપોર્ટ} (pāsporṭ)}{પાસપોર્ટ (pāsporṭ)}]{\gu{પાસપોર્ટ} (pāsporṭ) — passport}

\label{lesson:GU-C210-pasport}



\begin{quote}
Before the new one: say the Gujarati for to slide, then the Gujarati for airport.
\end{quote}

\subsection*{You'll want to know: \gu{પાસપોર્ટ}}
\textbf{\gu{પાસપોર્ટ}} — \emph{pāsporṭ} — \textquotedblleft{}passport\textquotedblright{}.

\begin{quote}
\textbf{\gu{મેં} \gu{પાસપોર્ટ} \gu{બતાવ્યો}.} — \emph{mẽ pāsporṭ batāvyo} — \textquotedblleft{}I showed my passport.\textquotedblright{}
\end{quote}

A masculine noun.

\begin{cousinweb}
\textbf{\gu{પાસપોર્ટ}} (\emph{pāsporṭ}) is a loan from English \emph{passport}.
\end{cousinweb}

\begin{grammarlens}[title={What you've built}]
One more word for everyday life at A2.
\end{grammarlens}

\subsection*{Your turn}
\begin{itemize}
\raggedright
  \item \emph{Say it:} \emph{pāsporṭ}
  \item \emph{Say it:} \emph{pāsporṭ}, once more
  \item \emph{Say it:} \emph{mẽ pāsporṭ batāvyo} — I showed my passport.
  \item \emph{Recall:} say the Gujarati for to slide, then the Gujarati for airport, then say \emph{pāsporṭ} again
\end{itemize}

\begin{tcolorbox}[breakable,colback=teal!4,colframe=teal!35!black,title={Before you move on}]
What does \gu{પાસપોર્ટ} mean? (\textquotedblleft{}Passport\textquotedblright{}.) Say it once more.
\end{tcolorbox}
\section[\texorpdfstring{\gu{પરવાનો} (parvāno)}{પરવાનો (parvāno)}]{\gu{પરવાનો} (parvāno) — permit, licence}

\label{lesson:GU-C210-parvano}



\begin{quote}
Before the new one: say the Gujarati for airport, then the Gujarati for passport.
\end{quote}

\subsection*{You'll want to know: \gu{પરવાનો}}
\textbf{\gu{પરવાનો}} — \emph{parvāno} — \textquotedblleft{}permit, licence\textquotedblright{}.

\begin{quote}
\textbf{\gu{ડ્રાઈવર} \gu{પાસે} \gu{પરવાનો} \gu{ન} \gu{હતો}.} — \emph{ḍraivar pāse parvāno na hato} — \textquotedblleft{}The driver had no licence.\textquotedblright{}
\end{quote}

\begin{cousinweb}
\textbf{\gu{પરવાનો}} (\emph{parvāno}) is a loan, usually traced to Persian \emph{parvāna}, \textquotedblleft{}warrant, licence\textquotedblright{}. CDIAL has no entry for it.
\end{cousinweb}

\begin{grammarlens}[title={What you've built}]
One more word for everyday life at A2.
\end{grammarlens}

\subsection*{Your turn}
\begin{itemize}
\raggedright
  \item \emph{Say it:} \emph{parvāno}
  \item \emph{Say it:} \emph{parvāno}, once more
  \item \emph{Say it:} \emph{ḍraivar pāse parvāno na hato} — The driver had no licence.
  \item \emph{Recall:} say the Gujarati for airport, then the Gujarati for passport, then say \emph{parvāno} again
\end{itemize}

\begin{tcolorbox}[breakable,colback=teal!4,colframe=teal!35!black,title={Before you move on}]
What does \gu{પરવાનો} mean? (\textquotedblleft{}Permit, licence\textquotedblright{}.) Say it once more.
\end{tcolorbox}
\section[\texorpdfstring{\gu{લાદવું} (lādvũ)}{લાદવું (lādvũ)}]{\gu{લાદવું} (lādvũ) — to load}

\label{lesson:GU-C210-ladvu}



\begin{quote}
Before the new one: say the Gujarati for passport, then the Gujarati for permit.
\end{quote}

\subsection*{You'll want to know: \gu{લાદવું}}
\textbf{\gu{લાદવું}} — \emph{lādvũ} — \textquotedblleft{}to load\textquotedblright{}.

\begin{quote}
\textbf{\gu{અમે} \gu{સામાન} \gu{ગાડીમાં} \gu{લાદ્યો}.} — \emph{ame sāmān gāḍīmā̃ lādyo} — \textquotedblleft{}We loaded the luggage into the car.\textquotedblright{}
\end{quote}

\begin{cousinweb}
\textbf{\gu{લાદવું}} (\emph{lādvũ}) is a tadbhava. It comes from Sanskrit \emph{lardayati}, \textquotedblleft{}loads\textquotedblright{}, and CDIAL glosses the Gujarati verb \textquotedblleft{}to heap up, load\textquotedblright{} (CDIAL 10966).
\end{cousinweb}

\begin{grammarlens}[title={What you've built}]
One more word for everyday life at A2.
\end{grammarlens}

\subsection*{Your turn}
\begin{itemize}
\raggedright
  \item \emph{Say it:} \emph{lādvũ}
  \item \emph{Say it:} \emph{lādvũ}, once more
  \item \emph{Say it:} \emph{ame sāmān gāḍīmā̃ lādyo} — We loaded the luggage into the car.
  \item \emph{Recall:} say the Gujarati for passport, then the Gujarati for permit, then say \emph{lādvũ} again
\end{itemize}

\begin{tcolorbox}[breakable,colback=teal!4,colframe=teal!35!black,title={Before you move on}]
What does \gu{લાદવું} mean? (\textquotedblleft{}To load\textquotedblright{}.) Say it once more.
\end{tcolorbox}
\section[\texorpdfstring{\gu{ઊંચકવું} (ū̃cakvũ)}{ઊંચકવું (ū̃cakvũ)}]{\gu{ઊંચકવું} (ū̃cakvũ) — to lift, to carry}

\label{lesson:GU-C210-uncakvu}



\begin{quote}
Before the new one: say the Gujarati for permit, then the Gujarati for to load.
\end{quote}

\subsection*{You'll want to know: \gu{ઊંચકવું}}
\textbf{\gu{ઊંચકવું}} — \emph{ū̃cakvũ} — \textquotedblleft{}to lift, to carry\textquotedblright{}.

\begin{quote}
\textbf{\gu{તેણે} \gu{ભારે} \gu{થેલી} \gu{ઊંચકી}.} — \emph{teṇe bhāre thelī ū̃cakī} — \textquotedblleft{}She lifted the heavy bag.\textquotedblright{}
\end{quote}

\begin{cousinweb}
\textbf{\gu{ઊંચકવું}} (\emph{ū̃cakvũ}) is a tadbhava. CDIAL lists it, glossed \textquotedblleft{}to raise\textquotedblright{}, under Sanskrit \emph{ucca}, \textquotedblleft{}tall, lofty\textquotedblright{}, built with the suffix \emph{-akka-} (CDIAL 1634).
\end{cousinweb}

\begin{grammarlens}[title={What you've built}]
One more word for everyday life at A2.
\end{grammarlens}

\subsection*{Your turn}
\begin{itemize}
\raggedright
  \item \emph{Say it:} \emph{ū̃cakvũ}
  \item \emph{Say it:} \emph{ū̃cakvũ}, once more
  \item \emph{Say it:} \emph{teṇe bhāre thelī ū̃cakī} — She lifted the heavy bag.
  \item \emph{Recall:} say the Gujarati for permit, then the Gujarati for to load, then say \emph{ū̃cakvũ} again
\end{itemize}

\begin{tcolorbox}[breakable,colback=teal!4,colframe=teal!35!black,title={Before you move on}]
What does \gu{ઊંચકવું} mean? (\textquotedblleft{}To lift, to carry\textquotedblright{}.) Say it once more.
\end{tcolorbox}
